\documentclass[10pt,a4paper]{amsart}
\usepackage[margin=19mm]{geometry}
\usepackage{amsmath,amssymb,mathtools}
\usepackage[scr=rsfs,cal=euler]{mathalfa}
\usepackage{xfrac}
\usepackage[unicode,hidelinks,bookmarks=false]{hyperref}
\newcommand{\ie}{i.e.\ }
\newcommand{\cf}{cf.\ }
\newcommand{\resp}{respectively\ }
\newcommand{\Subsection}{\subsection}
\providecommand{\nobreakdash}{\nobreak\hbox{-}}
\setlength{\emergencystretch}{2em}
\multlinegap0pt
\usepackage{bm}
\usepackage{bbm}
\usepackage{dsfont}
\usepackage{xspace}
\usepackage{cancel}
\usepackage{mathtools}
\usepackage{amsthm}
\makeatletter
\@ifundefined{theorem}{\newtheorem{theorem}{Theorem}}{}
\@ifundefined{lemma}{\newtheorem{lemma}[theorem]{Lemma}}{}
\makeatother
\title[Unbalanced Zarankiewicz problem for bipartite subdivisions w]{Unbalanced Zarankiewicz problem for bipartite subdivisions with applications to incidence geometry}
\author{Lili K\"odm\"on, Anqi Li, Ji Zeng}
\date{Source: arXiv:2412.10204v3 (CC BY 4.0)}
\begin{document}
\maketitle

\noindent\emph{Source and licence.} Unbalanced Zarankiewicz problem for bipartite subdivisions with applications to incidence geometry. Version arXiv:2412.10204v3 (CC BY 4.0), distributed under the Creative Commons Attribution 4.0 International licence, as recorded in the arXiv licence metadata for this exact version and shown on the abstract page https://arxiv.org/abs/2412.10204v3. Full rights notice: licenses/arxiv-2412.10204-CC-BY-4.0.txt.\par\smallskip

\noindent\textit{Editorial scope.} Counted unit: the Lemma whose source label is \texttt{biregular} in the article (source0.tex lines 92--99, source label \texttt{biregular}), delivered together with the article's own complete original proof of it (source0.tex lines 101--130, one \texttt{proof} environment of 30 lines). No mathematics is written, repaired or re-derived: statement and proof are the article's own text line for line. Required context retained uncounted: none. Every definition, display and label of those lines is retained unchanged. Omitted from this package and not redistributed: 1--91, 100--100, 131--469. Nothing inside a retained excerpt is omitted or altered: every retained body is delivered exactly as the article's lines are written, so no omission receipt of any kind accompanies this package. Citations: the retained excerpts carry no citation command at all, so no bibliography entry of the article is transcribed and no reference is redistributed. Reference adaptation: none was needed and none was made; every reference inside the delivered unit resolves inside it, so no unresolved reference or citation is produced. Printed slips kept literal: no printed word, symbol, hypothesis or number of the counted statement or of its proof was altered. Layout and class: the article's own class is \texttt{article} with packages including amssymb, amsmath, amsthm, graphicx, hyperref, geometry; the wrapper is the lane's pinned \texttt{amsart} editorial wrapper at 10pt on A4, which restates the theorem-like environments the retained text needs and adds the article's own macro definitions that the retained text needs (none), with extra wrapper packages (bm, bbm, dsfont, xspace, cancel, mathtools). The delivered numbering is the unit's own. Assets: no retained span contains a figure environment, a graphics-inclusion command, a PDF-page inclusion, a table or a TikZ picture; no image file is copied into this package, the package declares no asset dependency and no image file is redistributed with it. Licence: version arXiv:2412.10204v3 of this article is distributed under the Creative Commons Attribution 4.0 International licence, as recorded in the arXiv licence metadata for this exact version and shown on the abstract page https://arxiv.org/abs/2412.10204v3. A full rights notice is delivered at licenses/arxiv-2412.10204-CC-BY-4.0.txt. Characters: every retained excerpt is pure ASCII, so the delivered engine, XeLaTeX with its default Latin Modern text font, reports no missing character.
\begin{lemma}\label{biregular}
    For each integer $s \geq 2$, there is a sufficiently small positive constant $c_s$ such that any bipartite graph $G=(U \sqcup V,E)$ with $|V| \geq |U|^{2-1/s}$ and $|E|\geq \delta|V|$ ($1 \leq \delta \in \mathbb{R}$) contains a subgraph $\tilde{G} = (\tilde{U} \sqcup \tilde{V}, \tilde{E})$ with $|\tilde{U}| \geq |U|^{c_s}$ satisfying the following conditions: \begin{align*}
        &\makebox[0pt][l]{$\text{(I)}~|\tilde{V}| \geq |\tilde{U}|^{2-1/s},$}\phantom{\text{(III)}~c_s \Delta_{\tilde{G}}(\tilde{V}) \leq |\tilde{E}|/|\tilde{V}|,}
        \qquad\qquad \text{(II)}~|\tilde{E}|\geq c_s \delta|\tilde{V}|,\\
        &\text{(III)}~c_s \Delta_{\tilde{G}}(\tilde{V}) \leq |\tilde{E}|/|\tilde{V}|,
        \qquad\qquad \text{(IV)}~c_s \delta^{c_s} |\tilde{U}|\Delta_{\tilde{G}}(\tilde{U})^{1- 1/s} \leq |\tilde{E}|.
    \end{align*}
\end{lemma}

\begin{proof}
    We construct $\tilde{G}$ by two iterative procedures. In the first procedure, we let $G_0 = (U_0\sqcup V_0, E_0) = G$ and construct $G_i = (U_i\sqcup V_i, E_i)$ as follows. Given $G_i$ already constructed, we set $V_{i+1}$ to be the $|V_i|/16$ many vertices in $V_i$ with largest degrees. If less than half of $E_i$ are adjacent to $V_{i+1}$, we terminate the first procedure. Otherwise, by a simple averaging argument, we can find $U_{i+1} \subset U_i$ with $|U_{i+1}| = |U_i|/16^{\frac{s}{2s-1}}$ such that \begin{equation*}
        |E_{G_i}(U_{i+1}, V_{i+1})| \geq \frac{1}{2}|E_i|/16^{\frac{s}{2s-1}}.
    \end{equation*} Here, $E_{G_i}(U_{i+1}, V_{i+1})$ denotes the set of edges in $G_i$ between $U_{i+1}$ and $V_{i+1}$. Then we take $G_{i+1}$ as the subgraph of $G_i$ induced on $U_{i+1} \sqcup V_{i+1}$. Let $\ell$ be the index where the first procedure terminates at $G_{\ell}$. We have the following identities:\begin{equation*}
        |U_{\ell}| = |U_0|/16^{\frac{\ell s}{2s-1}},\quad |V_\ell| = |V_0|/16^\ell,\quad |E_\ell| \geq \frac{1}{2^\ell}|E_0|/16^{\frac{\ell s}{2s-1}}.
    \end{equation*} Combining them with $|U_\ell||V_\ell|\geq|E_\ell|$ and $|E_0| \geq |V_0|$, we can solve for \begin{equation}\label{eq:biregular_ell}
        \ell \leq \log(|U_0|)/3.
    \end{equation}

    Next, we let $\tilde{V} = V_\ell \setminus V_{\ell+1}$. Notice that $|\tilde{V}| = (15/16)|V_\ell|$ and $|E_{G_\ell}(U_\ell, \tilde{V})| \geq |E_\ell|/2$. By an averaging argument, we can find $U' \subset U_\ell$ with $|U'| = (15/16)^{\frac{s}{2s-1}}|U_\ell|$ such that the subgraph $G'\subset G_\ell$ induced on $U' \sqcup \tilde{V}$ satisfies $|E'| \geq (15/16)^{\frac{s}{2s-1}}|E_\ell|/2$, where $E'$ is the edge set of $G'$. Since $V_{\ell+1}$ is chosen among vertices with largest degrees, \begin{equation}\label{eq:biregular_B}
        \Delta_{G'}(\tilde{V}) \leq \frac{|E_{\ell}|}{|V_{\ell+1}|} \leq 30 \cdot (16/15)^{\frac{s}{2s-1}} \cdot \frac{|E'|}{|\tilde{V}|}.
    \end{equation}
    
    In the second procedure, we let $G'_0 = (U'_0 \sqcup \tilde{V}, E'_0) =G'$ and construct $G'_i = (U'_i \sqcup \tilde{V}, E'_i)$ as follows. Given $G'_i$ already constructed, we set $U'_{i+1}$ to be the $|U'_i|^{1-1/s}$ many vertices in $U'_i$ with largest degrees. If less than half of $E'_i$ are adjacent to $U'_{i+1}$, we terminate the second procedure. In this case, we let $\tilde{U} = U'_i \setminus U'_{i+1}$ and $\tilde{G}$ be the subgraph of $G'_i$ induced on $\tilde{U} \sqcup \tilde{V}$ and denote its edge set as $\tilde{E}$. Since $U'_{i+1}$ is chosen among vertices with largest degrees, we have \begin{equation}\label{eq:biregular_A1}
        \Delta_{\tilde{G}}(\tilde{U}) \leq \frac{|E'_i|}{|U'_{i+1}|} \leq \frac{2|\tilde{E}|}{|U'_{i+1}|} = \frac{2|\tilde{E}|}{|U'_i|^{1-1/s}} \leq \frac{2|\tilde{E}|}{|\tilde{U}|^{1-1/s}}.
    \end{equation} If more than half of $E'_i$ are adjacent to $U'_{i+1}$, we take $G'_{i+1}$ as the subgraph of $G'_i$ induced on $U'_{i+1}\sqcup \tilde{V}$. Suppose the second procedure has not terminated after $s\log(s)$ iterations, then we terminate it and define $\tilde{U} = U'_{s\log(s)}$ and $\tilde{G} = G'_{s\log(s)}$. In this case, we have \begin{equation}\label{eq:biregular_A2}
        |\tilde{U}| = |U'_0|^{(1-1/s)^{s\log(s)}} \leq |U'_0|^{1/s} \leq |\tilde{V}|^{1/s}.
    \end{equation}

    Now we check that $\tilde{G}$ satisfies the conditions we claimed. Using \eqref{eq:biregular_ell} from the first procedure, we can check that $|U_\ell| \geq |U_0|^{\frac{2s-3}{6s-3}}$. It is a consequence of the second procedure terminating in at most $s \log s $ iterations that $|\tilde{U}| \geq |U'_0|^{(1-1/s)^{s\log(s)}}$. Hence $|\tilde{U}| \geq |U|^{c_s}$ is satisfied for small enough $c_s$.

    It is easy to see that $|V_i| \geq |U_i|^{2-1/s}$ are guaranteed throughout the first procedure, and $|\tilde{V}| \geq |U'|^{2-1/s} \geq |\tilde{U}|^{2-1/s}$. So condition (I) is satisfied. We can inductively check that $|E_i| \geq \delta|V_i|$ in the first procedure. It is clear from the second procedure that $|\tilde{E}|$ and $|E_\ell|$ only differ by a multiplicative constant depending on $s$. This verifies condition (II) for small enough $c_s$. Similarly, we can verify condition (III) using \eqref{eq:biregular_B}.

    Finally, if the second procedure terminated with \eqref{eq:biregular_A1}, we can compute \begin{align*}
        |\tilde{U}|\Delta_{\tilde{G}}(\tilde{U})^{1 - 1/s} &\leq |\tilde{U}|\left( \frac{2|\tilde{E}|}{|\tilde{U}|^{1-1/s}} \right)^{1 - 1/s} = \left( |\tilde{U}|^{2-1/s} \right)^{1/s}\left( 2|\tilde{E}| \right)^{1 - 1/s} \leq |\tilde{V}|^{1/s}\left( 2|\tilde{E}| \right)^{1 - 1/s}\\
        &\leq \left( c_s^{-1}\delta^{-1}|\tilde{E}| \right)^{1/s}\left( 2|\tilde{E}| \right)^{1 - 1/s} = 2 \cdot \left( 2c_s\delta \right)^{-1/s}|\tilde{E}|.
    \end{align*} If the second procedure terminated with \eqref{eq:biregular_A2}, we have \begin{equation*}
         |\tilde{U}|\Delta_{\tilde{G}}(\tilde{U})^{1 - 1/s} \leq |\tilde{V}|^{1/s}|\tilde{V}|^{1-1/s} \leq |\tilde{V}| \leq c_s^{-1}\delta^{-1}|\tilde{E}|.
    \end{equation*} Hence we can verify condition (IV) for small enough $c_s$. Overall, the proof is finished.
\end{proof}

\end{document}
